%%
%% paper76_mandelbrot_mlc_en.tex
%% Autor: Michael Fuhrmann
%% Projekt: Specialist Mathematics Project
%% Thema: The Mandelbrot Set — Local Connectivity and the MLC Conjecture
%% Sprache: Englisch
%% Erstellt: 2026-03-12
%% Letzte Änderung: 2026-03-12
%% Build: 164
%%
%% Beschreibung:
%%   Dieses Paper behandelt die Mandelbrot-Menge M im Parameterraum
%%   der quadratischen Familie f_c(z) = z^2 + c, ihre globale Zusammenhängigkeit
%%   (Beweis von Douady-Hubbard 1982), die Vermutung der lokalen Zusammenhängigkeit
%%   (MLC), Resultate von Yoccoz, Thurston's topologisches Modell sowie
%%   offene Probleme in der holomorphen Dynamik.
%%

\documentclass[12pt,a4paper]{amsart}

%% Pakete für Mathematik, Layout, Encoding und Links
\usepackage{amsmath,amssymb,amsthm,mathtools,enumitem,booktabs,array,hyperref}
\usepackage[utf8]{inputenc}
\usepackage[T1]{fontenc}
\usepackage{geometry}
\geometry{margin=2.5cm}

%% Theorem-Umgebungen (englisch)
\newtheorem{theorem}{Theorem}[section]
\newtheorem{lemma}[theorem]{Lemma}
\newtheorem{corollary}[theorem]{Corollary}
\newtheorem{proposition}[theorem]{Proposition}
\theoremstyle{definition}
\newtheorem{definition}[theorem]{Definition}
\newtheorem{example}[theorem]{Example}
\theoremstyle{remark}
\newtheorem{remark}[theorem]{Remark}
\newtheorem{conjecture}[theorem]{Conjecture}

%% Mathematische Operatoren
\DeclareMathOperator{\diam}{diam}
\DeclareMathOperator{\dist}{dist}
\DeclareMathOperator{\Int}{Int}
\DeclareMathOperator{\Crit}{Crit}

%% Hyperref-Konfiguration
\hypersetup{
    colorlinks=true,
    linkcolor=blue,
    citecolor=blue,
    urlcolor=blue,
    pdftitle={The Mandelbrot Set: Local Connectivity and the MLC Conjecture},
    pdfauthor={Michael Fuhrmann}
}

%% -------------------------------------------------------------------------
%% Abstract VOR \maketitle
%% -------------------------------------------------------------------------
\begin{abstract}
We present a comprehensive treatment of the Mandelbrot set $\mathcal{M}$,
the canonical parameter space for the quadratic family $f_c(z) = z^2 + c$.
After reviewing the foundational notions of filled Julia sets and the
characterisation of $\mathcal{M}$ via non-escape of the critical orbit,
we give a detailed account of the Douady--Hubbard proof (1982) that
$\mathcal{M}$ is connected, via the Böttcher coordinate and the
conformal map $\Phi_{\mathcal{M}} \colon \widehat{\mathbb{C}} \setminus \mathcal{M}
\to \widehat{\mathbb{C}} \setminus \overline{\mathbb{D}}$.
We then formulate the \emph{MLC conjecture} — that $\mathcal{M}$ is
locally connected — and survey its profound implications: the density of
hyperbolic maps in the quadratic family, and Thurston's topological
characterisation of $\mathcal{M}$ via abstract Mandelbrot sets.
Yoccoz's theorem (MLC at all at most finitely renormalisable parameters),
the Fatou--Shishikura inequality, and the structure of satellite and
primitive components are discussed. We conclude with a list of open problems.
\end{abstract}

\title{The Mandelbrot Set: Local Connectivity and the MLC Conjecture}
\author{Michael Fuhrmann}
\address{Specialist Mathematics Project, Build 165}
\email{specialist-maths@project.local}
\date{2026-03-12}

\maketitle

%% -------------------------------------------------------------------------
\section{Introduction}
%% -------------------------------------------------------------------------

The \emph{Mandelbrot set} $\mathcal{M}$, introduced by Benoit Mandelbrot in
1980~\cite{Mandelbrot1980} and rigorously analysed by Douady and Hubbard
beginning in 1982~\cite{DouadyHubbard1982}, is the most celebrated object
in holomorphic dynamics. It parametrises the quadratic family
\[
    f_c(z) = z^2 + c, \quad c \in \mathbb{C},
\]
by marking those parameters for which the dynamical behaviour of $f_c$ is
``bounded.'' Despite its elementary definition, $\mathcal{M}$ encodes
an extraordinary wealth of complex-analytic structure. Its global
connectivity was established in the very first paper of Douady and Hubbard,
yet one of the most fundamental questions — whether $\mathcal{M}$ is
\emph{locally} connected — remains open after more than four decades.

This \emph{MLC conjecture} (``$\mathcal{M}$ is Locally Connected'') is not
merely a topological curiosity. It is equivalent to the \emph{density of
hyperbolicity} in the quadratic family: the statement that every $c \in \mathcal{M}$
can be approximated by hyperbolic parameters. Moreover, MLC would yield a
complete topological model of $\mathcal{M}$ via Thurston's theory of
laminations, reducing the combinatorial description of the set to a
purely combinatorial object: the \emph{Mandelbrot lamination}.

The present paper surveys the main results and open problems in this circle
of ideas. We assume familiarity with one complex variable and basic notions
of holomorphic dynamics; proofs of central theorems are sketched with
references to the original literature.

%% -------------------------------------------------------------------------
\section{The Quadratic Family and Julia Sets}
%% -------------------------------------------------------------------------

\subsection{Basic Definitions}

\begin{definition}[Filled Julia Set]
For $c \in \mathbb{C}$, the \emph{filled Julia set} of $f_c$ is
\[
    K_c = \{ z \in \mathbb{C} \mid \sup_{n \geq 0} |f_c^n(z)| < \infty \},
\]
and the \emph{Julia set} is $J_c = \partial K_c$.
\end{definition}

Since $f_c$ is a polynomial of degree 2, the complement
$\mathbb{C} \setminus K_c$ is the \emph{basin of infinity}
$\mathcal{A}_\infty(c)$. The Julia set $J_c$ is the boundary between
bounded and unbounded orbits. A fundamental dichotomy holds:

\begin{theorem}[Fatou--Julia Dichotomy]
For every $c \in \mathbb{C}$, either
\begin{enumerate}[label=\roman*)]
\item $K_c$ is connected and $J_c$ is a connected, perfect, nowhere
      dense subset of $\mathbb{C}$, or
\item $K_c$ is a Cantor set (totally disconnected).
\end{enumerate}
\end{theorem}

\begin{proof}[Proof sketch]
By the Riemann mapping theorem and the Böttcher coordinate, if $0 \in K_c$
(i.e.\ the critical point does not escape), then $K_c$ is full and connected.
If $0 \notin K_c$ (the critical orbit escapes), Böttcher's theorem
globalises to give a biholomorphism from $\mathbb{C} \setminus K_c$ to
$\mathbb{C} \setminus \overline{\mathbb{D}}$ with prescribed asymptotics,
showing $K_c$ is a Cantor set by Sullivan's no-wandering-domains theorem
(applicable in the polynomial case via direct arguments).
\end{proof}

\subsection{The Mandelbrot Set as Parameter Space}

\begin{definition}[Mandelbrot Set]
\[
    \mathcal{M} = \{ c \in \mathbb{C} \mid f_c^n(0) \not\to \infty \text{ as } n \to \infty \}
              = \{ c \in \mathbb{C} \mid K_c \text{ is connected} \}.
\]
\end{definition}

The equivalence in the definition is the content of the Fatou--Julia
dichotomy applied to the critical point $0$, which is the unique finite
critical point of $f_c$.

\begin{proposition}
$\mathcal{M}$ is compact, and $\mathcal{M} \subset \overline{D}(0, 2)$.
\end{proposition}

\begin{proof}
If $|c| > 2$, then $|f_c(0)| = |c| > 2$ and a standard estimate shows
$|f_c^n(0)| \to \infty$. Compactness follows from the continuity of
$c \mapsto f_c^n(0)$ and the definition as an intersection of closed sets.
\end{proof}

%% -------------------------------------------------------------------------
\section{The Böttcher Coordinate and Green's Function}
%% -------------------------------------------------------------------------

\begin{theorem}[Böttcher Coordinate]
For each $c \in \mathbb{C}$, there exists a unique conformal map
\[
    \phi_c \colon \mathcal{A}_\infty(c) \to \mathbb{C} \setminus \overline{\mathbb{D}}
\]
satisfying $\phi_c(f_c(z)) = \phi_c(z)^2$ and $\phi_c(z)/z \to 1$ as $z \to \infty$.
\end{theorem}

The \emph{Green's function} of $K_c$ is
\[
    G_c(z) = \log|\phi_c(z)| = \lim_{n\to\infty} \frac{1}{2^n} \log|f_c^n(z)|,
\]
which extends continuously to $\mathbb{C}$ by setting $G_c = 0$ on $K_c$.

\begin{definition}[External Rays]
For $\theta \in \mathbb{R}/\mathbb{Z}$, the \emph{external ray of angle
$\theta$} for $K_c$ is
\[
    \mathcal{R}_c(\theta) = \phi_c^{-1}\bigl(\{ r e^{2\pi i \theta} \mid r > 1 \}\bigr).
\]
An external ray \emph{lands} at a point $\gamma_c(\theta) \in J_c$ if
$\phi_c^{-1}(r e^{2\pi i \theta}) \to \gamma_c(\theta)$ as $r \to 1^+$.
\end{definition}

The dynamics of $f_c$ on rays is particularly clean: $f_c$ maps the ray of
angle $\theta$ to the ray of angle $2\theta \pmod 1$.

%% -------------------------------------------------------------------------
\section{Connectivity of the Mandelbrot Set}
%% -------------------------------------------------------------------------

\subsection{The Douady--Hubbard Conformal Map}

The global parameter-space analogue of the Böttcher coordinate is:

\begin{theorem}[Douady--Hubbard 1982~\cite{DouadyHubbard1982}]
\label{thm:DH-connected}
There exists a biholomorphism
\[
    \Phi_{\mathcal{M}} \colon \widehat{\mathbb{C}} \setminus \mathcal{M}
    \longrightarrow \widehat{\mathbb{C}} \setminus \overline{\mathbb{D}}
\]
satisfying $\Phi_{\mathcal{M}}(c)/c \to 1$ as $c \to \infty$, given
explicitly by $\Phi_{\mathcal{M}}(c) = \phi_c(c)$.
\end{theorem}

\begin{proof}[Proof sketch]
For $c \notin \mathcal{M}$, the critical value $c$ itself lies in the
basin of infinity of $f_c$, so $\phi_c(c)$ is well-defined. One must
verify that $c \mapsto \phi_c(c)$ is holomorphic in $c$ on
$\mathbb{C} \setminus \mathcal{M}$, injective, and has the correct
boundary behaviour. Holomorphicity follows from the implicit function
theorem applied to the Böttcher functional equation. Injectivity is the
key point: it uses the fact that $f_c$ has a unique critical point $0$
and a careful analysis of the dependence of the Böttcher coordinate on
the parameter. The map extends to a homeomorphism of the boundaries by
Carathéodory's theorem, since $\mathcal{M}$ is compact and full.
\end{proof}

\begin{corollary}[Connectivity of $\mathcal{M}$]
The Mandelbrot set $\mathcal{M}$ is connected.
\end{corollary}

\begin{proof}
The existence of a conformal map from the complement of $\mathcal{M}$ to
the complement of $\overline{\mathbb{D}}$ implies, by Riemann's theory
of simply-connected domains, that $\mathcal{M}$ is connected and
full (i.e.\ $\mathbb{C} \setminus \mathcal{M}$ is connected).
\end{proof}

\subsection{External Rays of the Mandelbrot Set}

The conformal map $\Phi_{\mathcal{M}}$ defines external rays for $\mathcal{M}$:
\[
    \mathcal{R}_{\mathcal{M}}(\theta) = \Phi_{\mathcal{M}}^{-1}\bigl(\{ r e^{2\pi i \theta} \mid r > 1 \}\bigr).
\]
A ray $\mathcal{R}_{\mathcal{M}}(\theta)$ \emph{lands} at $c_0 \in \partial \mathcal{M}$
if $\Phi_{\mathcal{M}}^{-1}(r e^{2\pi i\theta}) \to c_0$ as $r \to 1^+$.

\begin{theorem}[Landing of Rational Rays]
Every external ray $\mathcal{R}_{\mathcal{M}}(\theta)$ with $\theta \in \mathbb{Q}/\mathbb{Z}$
lands at a Misiurewicz point or at the root of a hyperbolic component.
\end{theorem}

This is a consequence of the Douady--Hubbard theory and the classification
of periodic points on $\partial\mathcal{M}$.

%% -------------------------------------------------------------------------
\section{Hyperbolic Components and Internal Structure}
%% -------------------------------------------------------------------------

\begin{definition}[Hyperbolic Map]
$f_c$ is \emph{hyperbolic} if every critical orbit converges to an
attracting cycle. The \emph{main cardioid} $\mathcal{H}_0$ and the
\emph{period-2 bulb} $\mathcal{H}_2$ are the largest hyperbolic
components of $\mathcal{M}$.
\end{definition}

\begin{theorem}[Parametrisation of Hyperbolic Components]
Every hyperbolic component $\mathcal{H}$ of $\mathcal{M}$ is simply
connected, and there exists a biholomorphism
$\lambda_{\mathcal{H}} \colon \mathcal{H} \to \mathbb{D}$,
the \emph{multiplier map}, sending $c$ to the multiplier $\lambda$ of the
attracting cycle of $f_c$.
\end{theorem}

\begin{proof}
For the main cardioid $\mathcal{H}_0$, $f_c$ has an attracting fixed point
$z_c$ with multiplier $\lambda = f_c'(z_c) = 2z_c$. Solving $z^2+c=z$ gives
$z_c = (1 - \sqrt{1-4c})/2$, so the multiplier map is
$c = \lambda/2 - \lambda^2/4$, a biholomorphism from $\mathcal{H}_0$ to
$\mathbb{D}$. The general case follows by the same implicit-function-theorem
argument applied to the multiplier of the attracting cycle.
\end{proof}

%% -------------------------------------------------------------------------
\section{The MLC Conjecture}
%% -------------------------------------------------------------------------

\subsection{Statement and Equivalent Formulations}

\begin{conjecture}[MLC — Mandelbrot Locally Connected]
\label{conj:MLC}
The Mandelbrot set $\mathcal{M}$ is locally connected, i.e.\ for every
$c_0 \in \mathcal{M}$ and every neighbourhood $U$ of $c_0$, there
exists a connected neighbourhood $V \subseteq U$ of $c_0$ with
$V \cap \mathcal{M}$ connected.
\end{conjecture}

Local connectivity of $\mathcal{M}$ is equivalent to the following:

\begin{theorem}[MLC Equivalences]
The following statements are equivalent:
\begin{enumerate}[label=\roman*)]
\item $\mathcal{M}$ is locally connected (MLC).
\item The conformal map $\Phi_{\mathcal{M}}^{-1} \colon \mathbb{D}^c \to \mathcal{M}^c$
      extends continuously to the boundary: $\Phi_{\mathcal{M}}^{-1}$ extends
      to a continuous surjection $\overline{\mathbb{D}^c} \to \overline{\mathcal{M}^c}$,
      hence a Carathéodory loop $\mathbb{S}^1 \to \partial\mathcal{M}$.
\item Every external ray of $\mathcal{M}$ lands, and the landing map
      $\theta \mapsto \gamma_{\mathcal{M}}(\theta)$ is continuous.
\item The combinatorial model $\mathcal{M}_{\mathrm{comb}}$ (Thurston's
      abstract Mandelbrot set) is homeomorphic to $\mathcal{M}$.
\end{enumerate}
\end{theorem}

\begin{proof}[Proof of equivalences (sketch)]
The equivalence of (i), (ii), (iii) is standard in complex analysis:
for a compact connected set $K$ with connected complement, local connectivity
is equivalent to the Carathéodory extension property
(Carathéodory's theorem applied to the Riemann map of $\mathbb{C} \setminus K$).
The equivalence with (iv) is the content of Thurston's work on
topological polynomials and laminations~\cite{Thurston1985}.
\end{proof}

\subsection{Density of Hyperbolicity}

\begin{conjecture}[Density of Hyperbolicity]
\label{conj:hyp-density}
Hyperbolic maps are dense in the quadratic family: for every
$c_0 \in \mathcal{M}$ and every $\varepsilon > 0$, there exists
$c \in \mathcal{M}$ with $|c - c_0| < \varepsilon$ and $f_c$ hyperbolic.
\end{conjecture}

\begin{theorem}[MLC implies Density of Hyperbolicity]
Conjecture~\ref{conj:MLC} implies Conjecture~\ref{conj:hyp-density}.
\end{theorem}

\begin{proof}[Proof sketch]
Assuming MLC, the impression of every external ray of $\mathcal{M}$ is a
single point. The boundaries of hyperbolic components are dense in
$\partial \mathcal{M}$ (a consequence of the parametrisation of
hyperbolic components and the rational-ray landing theorem). Every point
in $\partial \mathcal{M}$ is then a limit of roots of hyperbolic components.
Interior points of $\mathcal{M}$ that are not in any hyperbolic component
would, assuming MLC, form a nowhere-dense set (since the Mandelbrot set
has no ``wandering'' Fatou components, by Sullivan's theorem). Together
this gives density of hyperbolic parameters.
\end{proof}

%% -------------------------------------------------------------------------
\section{Thurston's Topological Model}
%% -------------------------------------------------------------------------

\subsection{Laminations}

\begin{definition}[Quadratic Invariant Lamination]
A \emph{quadratic invariant lamination} (QIL) is a closed collection
$\mathcal{L}$ of chords of $\overline{\mathbb{D}}$ (``leaves'') such that:
\begin{enumerate}[label=\roman*)]
\item Leaves do not cross in $\mathbb{D}$.
\item $\mathcal{L}$ is invariant under $z \mapsto z^2$ on $\partial \mathbb{D}$.
\item The forward image of each leaf is a leaf.
\end{enumerate}
The \emph{Mandelbrot lamination} $\mathcal{L}_{\mathcal{M}}$ is the
QIL encoding the combinatorics of $\mathcal{M}$.
\end{definition}

\begin{theorem}[Thurston~\cite{Thurston1985}]
The quotient $\partial \mathbb{D} / \sim_{\mathcal{L}_{\mathcal{M}}}$,
where $\theta_1 \sim \theta_2$ iff $\theta_1$ and $\theta_2$ are
endpoints of a leaf of $\mathcal{L}_{\mathcal{M}}$ or in the same
``gap,'' is a topological model $\mathcal{M}_{\mathrm{comb}}$
homeomorphic to a locally connected copy of the Mandelbrot set.
If MLC holds, then $\mathcal{M}_{\mathrm{comb}} \cong \mathcal{M}$.
\end{theorem}

%% -------------------------------------------------------------------------
\section{Yoccoz's Theorem}
%% -------------------------------------------------------------------------

The most significant partial result towards MLC is due to Jean-Christophe
Yoccoz (1990, unpublished; see~\cite{HubbardSchleicher1994,Milnor2000}):

\begin{theorem}[Yoccoz~\cite{HubbardSchleicher1994}]
\label{thm:Yoccoz}
Let $c_0 \in \mathcal{M}$ be a parameter that is neither in the closure
of a hyperbolic component nor infinitely renormalisable. Then $\mathcal{M}$
is locally connected at $c_0$.
\end{theorem}

\begin{proof}[Proof strategy]
Yoccoz introduces the \emph{Yoccoz puzzle}, a sequence of nested
topological discs $P^0 \supset P^1 \supset \cdots$ in the dynamical plane
of $f_{c_0}$, cut along external rays landing at preimages of the $\alpha$
fixed point. The \emph{Yoccoz inequality}
\[
    \mathrm{mod}(P^{n+1} \setminus \overline{P^{n+2}}) \geq
    \mathrm{mod}(P^n \setminus \overline{P^{n+1}})
\]
shows that the moduli are bounded below, which forces the puzzle pieces
to shrink to a point. The corresponding pieces in the \emph{parameter puzzle}
then also shrink, proving local connectivity of $\mathcal{M}$ at $c_0$.
The finiteness of renormalisation is precisely the hypothesis that prevents
the puzzle from degenerating.
\end{proof}

\begin{corollary}
MLC holds for all Misiurewicz points and for all $c_0$ with non-recurrent
critical point (i.e.\ $0 \notin \omega(0)$ under $f_{c_0}$).
\end{corollary}

%% -------------------------------------------------------------------------
\section{Renormalisation and Infinitely Renormalisable Parameters}
%% -------------------------------------------------------------------------

\begin{definition}[Renormalisable Map]
$f_c$ is \emph{$n$-renormalisable} if there exists a topological disc
$U \ni 0$ such that $f_c^n \colon U \to f_c^n(U)$ is a polynomial-like
map of degree 2 with connected Julia set. $f_c$ is \emph{infinitely
renormalisable} if it is $n$-renormalisable for infinitely many $n$.
\end{definition}

Yoccoz's theorem leaves open the case of infinitely renormalisable
parameters. For these, no general proof of MLC is known. Partial results
include:

\begin{theorem}[Lyubich 1997~\cite{Lyubich1997}]
For parameters of \emph{bounded combinatorial type} (i.e.\ the renormalisation
periods are bounded), MLC holds.
\end{theorem}

\begin{theorem}[Kahn--Lyubich 2009~\cite{KahnLyubich2009}]
The Julia set $J_c$ is locally connected for every $c \in \mathcal{M}$
that is at most finitely renormalisable.
\end{theorem}

The infinitely renormalisable Feigenbaum parameter $c_{\mathrm{Feig}} \approx -1.4011552$
is the primary obstruction: MLC at this parameter would follow from
bounds on the conformal moduli of the Feigenbaum fixed point, a problem
still unresolved in full generality.

%% -------------------------------------------------------------------------
\section{The Fatou--Shishikura Inequality}
%% -------------------------------------------------------------------------

\begin{theorem}[Fatou--Shishikura Inequality~\cite{Shishikura1987}]
For any rational map $f$ of degree $d \geq 2$, the number of
non-repelling cycles satisfies
\[
    \#\{\text{non-repelling cycles}\} \leq 2d - 2.
\]
\end{theorem}

\begin{proof}[Proof for degree 2]
For $f_c$ with $d=2$, we have $2d-2 = 2$ critical points counted with
multiplicity (one finite critical point $0$ of multiplicity 1, and one at
$\infty$). Shishikura's proof uses quasiconformal deformation theory: one
shows that each non-repelling cycle ``absorbs'' at least one critical orbit
via a holomorphic motion argument, and distinct non-repelling cycles absorb
distinct critical orbits. For $f_c$ this gives at most 1 finite non-repelling
cycle, consistent with the structure of $\mathcal{M}$.
\end{proof}

\begin{corollary}
Every $f_c$ has at most one attracting or indifferent cycle. In particular,
if $c \in \mathcal{M}$ and $f_c$ has an attracting cycle, then $K_c$ is
a Jordan curve and $f_c$ is hyperbolic.
\end{corollary}

%% -------------------------------------------------------------------------
\section{Structural Results on \texorpdfstring{$\partial\mathcal{M}$}{∂M}}
%% -------------------------------------------------------------------------

\subsection{Misiurewicz Points}

\begin{definition}[Misiurewicz Point]
$c_0 \in \partial \mathcal{M}$ is a \emph{Misiurewicz point} if $0$ is
strictly pre-periodic under $f_{c_0}$: $f_{c_0}^{m+k}(0) = f_{c_0}^m(0)$
for some $m \geq 1$, $k \geq 1$.
\end{definition}

\begin{theorem}
Misiurewicz points are dense in $\partial \mathcal{M}$.
\end{theorem}

\subsection{Parabolic Implosion}

At parabolic parameters (where $f_c$ has a parabolic cycle with multiplier
$e^{2\pi i p/q}$), the Julia set $J_c$ undergoes a dramatic change:
infinitely many tiny copies of $J_c$ implode into the Fatou component,
a phenomenon known as \emph{parabolic implosion}~\cite{LavaursRousseau1993}.
This gives the characteristic ``antenna'' structure near parabolic points
in $\partial\mathcal{M}$.

\subsection{Hausdorff Dimension of $\partial \mathcal{M}$}

\begin{theorem}[Shishikura 1998~\cite{Shishikura1998}]
The Hausdorff dimension of the boundary of the Mandelbrot set equals 2:
\[
    \dim_H(\partial \mathcal{M}) = 2.
\]
\end{theorem}

%% -------------------------------------------------------------------------
\section{Open Problems}
%% -------------------------------------------------------------------------

We collect the most important open problems in the area.

\begin{conjecture}[MLC, Douady~\cite{DouadyHubbard1982}]
\label{conj:MLC2}
The Mandelbrot set is locally connected.
\end{conjecture}

\begin{conjecture}[Density of Hyperbolicity]
\label{conj:DHyp}
Hyperbolic parameters are dense in $\mathcal{M}$.
\end{conjecture}

\begin{conjecture}[MLC at the Feigenbaum Point]
The Mandelbrot set is locally connected at the Feigenbaum parameter
$c_{\mathrm{Feig}} \approx -1.4011552$. Equivalently, the conformal
moduli of the Feigenbaum polynomial-like restriction grow without bound.
\end{conjecture}

\begin{conjecture}[Siegel Discs on $\partial\mathcal{M}$]
For every $c$ on the boundary of a Siegel disc component, $J_c$ is a
Jordan curve.
\end{conjecture}

\begin{remark}
Conjecture~\ref{conj:MLC2} implies Conjecture~\ref{conj:DHyp}, but the
converse is not known. Avila--Lyubich--de Melo have made substantial
progress on renormalisation, but MLC at infinitely renormalisable points
of unbounded type remains the central open problem in holomorphic dynamics.
\end{remark}

%% -------------------------------------------------------------------------
\section{Conclusion}
%% -------------------------------------------------------------------------

The Mandelbrot set is simultaneously a parameter space, a combinatorial
encyclopedia of all quadratic dynamics, and a concrete object revealing the
deepest features of holomorphic iteration. Its connectivity
(Douady--Hubbard) is fully established; its local connectivity (MLC) is the
grand open problem of the field. Yoccoz's puzzle technique resolved MLC for
all at-most-finitely-renormalisable parameters, reducing the problem to
the infinitely renormalisable case. The equivalence with density of
hyperbolicity shows that MLC is not merely a topological statement but
a structural assertion about the entire quadratic family. Its resolution
would complete our understanding of quadratic dynamics in a manner parallel
to what the Riemann Hypothesis would achieve for the distribution of primes.

\begin{thebibliography}{99}

\bibitem{Mandelbrot1980}
B.~Mandelbrot,
\textit{Fractal aspects of the iteration of $z \mapsto \lambda z(1-z)$ for
complex $\lambda$ and $z$},
Ann.\ N.Y.\ Acad.\ Sci.\ \textbf{357} (1980), 249--259.

\bibitem{DouadyHubbard1982}
A.~Douady, J.~H.~Hubbard,
\textit{Itération des polynômes quadratiques complexes},
C.\ R.\ Acad.\ Sci.\ Paris \textbf{294} (1982), 123--126.

\bibitem{DouadyHubbard1984}
A.~Douady, J.~H.~Hubbard,
\textit{Étude dynamique des polynômes complexes, I, II},
Publications Mathématiques d'Orsay \textbf{84-02}, \textbf{85-04},
Université de Paris-Sud, Orsay, 1984--1985.

\bibitem{Thurston1985}
W.~P.~Thurston,
\textit{On the combinatorics and dynamics of iterated rational maps},
Preprint, Princeton University, 1985. Published in: D.~Schleicher (ed.),
\textit{Complex Dynamics: Families and Friends}, A.~K.~Peters, 2009, 3--137.

\bibitem{Milnor2000}
J.~Milnor,
\textit{Local connectivity of Julia sets: expository lectures},
in: \textit{The Mandelbrot Set, Theme and Variations},
London Math.\ Soc.\ Lecture Note Ser.\ \textbf{274},
Cambridge Univ.\ Press, 2000, 67--116.

\bibitem{HubbardSchleicher1994}
J.~H.~Hubbard, D.~Schleicher,
\textit{The spider algorithm},
in: \textit{Complex Dynamical Systems: The Mathematics Behind the
Mandelbrot and Julia Sets}, Proc.\ Sympos.\ Appl.\ Math.\ \textbf{49},
Amer.\ Math.\ Soc., 1994, 155--180.

\bibitem{Lyubich1997}
M.~Lyubich,
\textit{Dynamics of quadratic polynomials, I–II},
Acta Math.\ \textbf{178} (1997), 185--247 and 247--297.

\bibitem{KahnLyubich2009}
J.~Kahn, M.~Lyubich,
\textit{The quasi-additivity law in conformal geometry},
Ann.\ of Math.\ \textbf{169} (2009), 561--593.

\bibitem{Shishikura1987}
M.~Shishikura,
\textit{On the quasiconformal surgery of rational functions},
Ann.\ Sci.\ École Norm.\ Sup.\ \textbf{20} (1987), 1--29.

\bibitem{Shishikura1998}
M.~Shishikura,
\textit{The Hausdorff dimension of the boundary of the Mandelbrot set
and Julia sets},
Ann.\ of Math.\ \textbf{147} (1998), 225--267.

\bibitem{LavaursRousseau1993}
P.~Lavaurs,
\textit{Systèmes dynamiques holomorphes: explosion de points périodiques
paraboliques},
Thèse, Université Paris-Sud, Orsay, 1989.

\bibitem{AvilaLyubich2008}
A.~Avila, M.~Lyubich,
\textit{Feigenbaum Julia sets of positive area},
Preprint arXiv:math/0602458, 2006.

\end{thebibliography}

\end{document}
